% Generated draft: requires mathematical and source review.
\documentclass[11pt]{article}
\usepackage[margin=25mm]{geometry}
\usepackage{fontspec}
\setmainfont{Latin Modern Roman}
\usepackage{amsmath,amssymb,mathtools,unicode-math}
\setmathfont{Latin Modern Math}
\usepackage{enumitem,needspace,xurl,hyperref}
\setlength{\parindent}{0pt}
\setlength{\parskip}{5pt}
\setlength{\emergencystretch}{3em}
\newcommand{\N}{\mathbb{N}}
\newcommand{\Z}{\mathbb{Z}}
\newcommand{\Q}{\mathbb{Q}}
\newcommand{\R}{\mathbb{R}}
\newcommand{\C}{\mathbb{C}}
\newcommand{\F}{\mathbb{F}}
\newcommand{\A}{\mathbb{A}}
\newcommand{\PP}{\mathbb{P}}
\newcommand{\E}{\mathbb{E}}
\newcommand{\eps}{\varepsilon}
\newcommand{\dd}{\,\mathrm{d}}
\begin{document}
\section*{UnsolvedMath ID: 30003774}
\textbf{Source problem code:} catalogue entry
\subsection{Definitions and mathematical statement}
Let \(p\in\R\) satisfy \(1<p<5\). The source concerns the one-dimensional nonlinear Schrödinger equation
\[
 i\partial_s U+\partial_y^2U=|U|^{p-1}U,
 \qquad U(0,\cdot)=U_0.
\]
Here \(s\) denotes time, \(y\) denotes the spatial variable, and \(U_0\) is described as Gaussian random initial data of low regularity. The spatial domain, boundary conditions, Gaussian law, and precise regularity space are not specified in the supplied record.

To express the probabilistic qualification, write \(U_0\) as a random datum on a probability space \((\Omega,\mathcal F,\mathbb P)\). The phrase ``almost-sure global solution'' means that the intended solution exists globally in time for \(\mathbb P\)-almost every realization of the datum, in an intended solution class that the record does not identify. This describes the premise of the recorded question; it is not an independently established existence or uniqueness assertion here.

\textbf{Recorded problem.} For each exponent \(1<p<5\), determine the long-time behavior of the almost-sure global solution \(U\) associated with the intended Gaussian low-regularity initial data. The intended probability law is unspecified; this question must not be read as quantifying over every Gaussian law. The record does not specify a limiting profile, a topology of convergence, a definition of scattering, or whether one or both time directions are intended. Consequently, it supplies a research question but does not determine a unique precise asymptotic proposition.

\textbf{Supplied status qualification.} The record labels the problem partially solved. Its dated assessment attributes global well-posedness for \(p>1\) and almost-sure scattering for \(p>3\) to Burq--Thomann, while stating that the broad long-time question is not resolved in every regime or for arbitrary data. These are supplied literature claims, not verified facts in this draft. Their exact data classes and other hypotheses are absent, so they cannot be applied here as unconditional assertions about the unspecified Gaussian data. The mention of arbitrary data does not enlarge the recorded random-data problem.
\subsection{Short English statement}
For \(1<p<5\), describe the long-time behavior of the one-dimensional nonlinear Schrödinger solution arising almost surely from the intended Gaussian initial data of low regularity. The domain, random-data law, regularity, solution class, and desired asymptotic conclusion remain unspecified. The supplied partially solved label and literature claims require independent review.
\subsection{Sources}
\begin{itemize}
\item Local UnsolvedMath registry: supplied complete source record for permanent ID \(30003774\), source problem code \texttt{catalogue entry}, in the Oberwolfach Reports Open Problems collection. No individual proposer is named. The record gives proposed year 2018 and labels its statement \texttt{corrected\_verified}; its verification note describes repairs to notation, typography, and sentence structure while preserving mathematical content. This is supplied editorial metadata, not independent mathematical verification.
\item Supplied source citation: \emph{Mini-Workshop: Gibbs Measures for Nonlinear Dispersive Equations (2019)}, Oberwolfach Reports, pp.\ 1081--1116, DOI \texttt{10.4171/owr/2018/18}. \url{https://doi.org/10.4171/owr/2018/18}. The background describes Oberwolfach Reports as documenting mathematical workshops and their open-problem discussions. The literature-source label is ``OWR 2018 report''; the supplied citation's 2019 designation and the record's 2018 provenance are retained without reconciliation.
\item Supplied literature reference: Burq--Thomann, ``almost-sure scattering,'' DOI \texttt{10.1090/memo/1480}. \url{https://doi.org/10.1090/memo/1480}. The record attributes global well-posedness for \(p>1\) and almost-sure scattering for \(p>3\) to this reference, without supplying the complete theorem hypotheses.
\item Literature assessment checked at \texttt{2026-08-21}: the supplied curation labels the problem partially solved and says the broad long-time question remains unresolved in some regimes or for arbitrary data. The background explicitly describes this as dated literature triage derived from the supplied Oberwolfach curation and calls for independent verification of time-sensitive status and cited claims. This date and assessment are supplied dated curation, not independently verified; no source visit is claimed.
\end{itemize}
\end{document}
